\section{Dataset and Data Engineering}
\label{sec:dataset}

\subsection{Data Sources}

The dataset is assembled from three heterogeneous sources:

\begin{enumerate}[leftmargin=1.5em]
  \item \textbf{Laboratory captures.}
    Images taken inside experimental cages at Shanghai Jiao Tong University
    School of Medicine, annotated in-house with VOC XML labels.
    These cover diverse lighting conditions and cage configurations.

  \item \textbf{Open-source datasets.}
    Publicly available rodent image collections, re-annotated and mapped
    to our two-class schema.

  \item \textbf{Web-crawled images.}
    Photographs retrieved via keyword search, manually inspected and
    cleaned to remove near-duplicates and mislabelled samples.
\end{enumerate}

\subsection{Preprocessing and Annotation Unification}

All images are normalised to Pascal VOC format~\citep{everingham2010pascal}
with a unified two-class label vocabulary:
\texttt{mouse} (experimental mouse, \textit{Mus musculus}) and
\texttt{other} (any non-mouse foreground object).
Preprocessing steps include:

\begin{itemize}[leftmargin=1.5em]
  \item Label harmonisation: merging synonymous class names (\textit{rat},
    \textit{mice}, \textit{rodent}~$\to$~\texttt{mouse}).
  \item Near-duplicate removal using perceptual hashing with a Hamming
    distance threshold of 4.
  \item Bounding-box quality filtering: boxes with aspect ratio $>$10 or
    area $<$ 400\,px\textsuperscript{2} are discarded.
\end{itemize}

\subsection{Dataset Statistics}

\begin{table}[H]
\centering
\caption{Dataset summary after preprocessing.}
\label{tab:dataset}
\begin{tabular}{lrr}
\toprule
\textbf{Split} & \textbf{Images} & \textbf{Fraction} \\
\midrule
Training (full) & 5{,}828 & 79.9\% \\
Validation      & 1{,}458 & 20.1\% \\
\midrule
\textbf{Total}  & \textbf{7{,}286} & 100\% \\
\bottomrule
\end{tabular}
\end{table}

\subsection{One-Third Subset for Ablation}
\label{sec:dataset:subset}

To quantify the effect of dataset scale on YOLOv3, we generate a
one-third training subset of 1,942 images using a script that performs
stratified random sampling at the file level
(\texttt{scripts/prepare\_1of3\_dataset.py}).
The \emph{validation set is shared} between the full and subset
experiments, ensuring that mAP values are directly comparable.
This design isolates training-set size as the sole variable between
experiment groups Y1/Y2 (1,942 train images) and Y3/Y4 (5,828 train images).

\begin{figure}[H]
\centering
\begin{tikzpicture}[font=\small]
  % Full dataset box
  \node[draw, fill=picoBlue!15, rounded corners, minimum width=5cm,
        minimum height=1.2cm, align=center] (full)
        at (0,0) {\textbf{Full Dataset}\\7,286 images};

  % Split arrow
  \draw[->, thick] (full.south) -- ++(0,-0.6)
    node[below, xshift=-2.4cm] {}
    node[below, xshift=2.4cm] {};

  \node[draw, fill=picoBlue!25, rounded corners, minimum width=2.8cm,
        minimum height=1cm, align=center] (train)
        at (-2.5,-2.0) {\textbf{Train}\\5,828};
  \node[draw, fill=accentGreen!25, rounded corners, minimum width=2.8cm,
        minimum height=1cm, align=center] (val)
        at (2.5,-2.0) {\textbf{Val}\\1,458};

  \draw[->, thick] (0,-0.6) -- (train.north);
  \draw[->, thick] (0,-0.6) -- (val.north);

  % 1/3 subset
  \node[draw, fill=yoloOrange!20, rounded corners, minimum width=2.8cm,
        minimum height=1cm, align=center, dashed] (subset)
        at (-2.5,-3.6) {\textbf{1/3 Subset}\\1,942};
  \draw[->, dashed, thick] (train.south) -- (subset.north)
    node[midway, right, xshift=2pt] {\scriptsize stratified sample};

  % Shared val annotation
  \draw[<->, thick, accentGreen!70!black] (subset.east) -- (val.south)
    node[midway, below, sloped, xshift=6pt] {\scriptsize shared val};

  % Labels
  \node[above right=0.1cm of full, font=\scriptsize\itshape]
    {All sources unified to VOC format};
\end{tikzpicture}
\caption{Dataset split design.  The validation set is shared between
full and one-third-subset experiments to enable direct mAP comparison.}
\label{fig:dataset_split}
\end{figure}
